\section*{Hoofdstuk \ref{ch:g2:caring-for-nature} --- Zorgen voor de natuur}

\begin{solution}{exo:g2:caring-for-nature:1}
De omgeving waarvan een \omterm{def:g1:living-or-not:living}{levend wezen} afhangt. Bijvoorbeeld: de lucht
die we inademen, het water dat we drinken, het voedsel dat we eten ---
allemaal uit de \omterm{def:g2:caring-for-nature:environment}{natuur}.
\end{solution}

\begin{solution}{exo:g2:caring-for-nature:2}
Bijvoorbeeld: de kevers onder de bast, de pissebedden onder de stam,
het mos erbovenop, de egel die in de holte overwintert.
\end{solution}

\begin{solution}{exo:g2:caring-for-nature:3}
Bijvoorbeeld: neem al je afval mee naar huis; kijk in plaats van te
verzamelen, en pluk nooit een hele pol leeg; leg elke omgedraaide steen
of stam terug; blijf op de paden bij kwetsbare plekken; kijk stil van
een afstand naar \omterm{def:g1:animals-around-us:animal}{dieren} en laat hun jongen met rust.
\end{solution}

\begin{solution}{exo:g2:caring-for-nature:4}
Omdat de \omterm{def:g2:caring-for-nature:environment}{natuur} geen vuilniswagen heeft: wat we achterlaten blijft ---
een plastic dop heel veel jaren, lang genoeg om \omterm{def:g1:animals-around-us:animal}{dieren} te vangen of te
vergiftigen en er nog te liggen als de kinderen van nu volwassen zijn.
\end{solution}

\begin{solution}{exo:g2:caring-for-nature:5}
Tot de laatste geplukt kan de pol geen \omterm{def:g2:from-seed-to-plant:seed}{zaad} meer maken --- dus volgend
jaar geen madeliefjes daar --- en de bijen raken de \omterm{def:g1:plants-around-us:parts}{bloemen} kwijt
waarvan ze eten.
\end{solution}

\begin{solution}{exo:g2:caring-for-nature:6}
Ze verteren langzaam tot donkere, kruimelige aarde. Wormen, pissebedden
en levende wezens die te klein zijn om te zien doen het werk; aan het
eind krijgen we compost --- goede aarde die de \omterm{def:g1:plants-around-us:plant}{planten} van volgend jaar voedt.
\end{solution}

\begin{solution}{exo:g2:caring-for-nature:7}
Bijvoorbeeld: een schoteltje water tijdens een hittegolf, een wild,
ongemaaid hoekje in de tuin, zaadjes voor vogels in de kou, een klein
vijvertje graven.
\end{solution}

\begin{solution}{exo:g2:caring-for-nature:8}
Open antwoord. Plastic verpakkingen en flessen winnen meestal. Elke
soort had met zijn eigenaar mee naar huis gemoeten --- en daarna in de
juiste sorteerbak.
\end{solution}

\begin{solution}{exo:g2:caring-for-nature:9}
Laat de \omterm{ex:g2:animals-grow-up:frog}{donderpadden} in de vijver: die heeft alles wat ze nodig hebben
--- vijverwater, voedsel, ruimte om poten te krijgen --- en een pot
niet. Als Emma wil helpen, kan ze het voorjaar lang naar ze kijken, de
oever onvertrapt houden, en misschien meehelpen een nieuw
vijverleefgebied te maken zoals dat van de school.
\end{solution}

\begin{solution}{exo:g2:caring-for-nature:10}
Een miljoen ``nietsjes'' de verkeerde kant op is een wei bedolven onder
doppen --- kleine schade vermenigvuldigt zich. De goede kant op tellen
een miljoen kleine zorgzaamheden --- afval mee naar huis, kranen dicht,
afval gesorteerd --- net zo zeker op: de \omterm{def:g2:caring-for-nature:environment}{natuur} wordt vooral verzorgd
in keuzes die precies zo klein zijn.
\end{solution}
